% TOP_PROBLEM: {"id": 2884, "title": "Kirby Problem 4.8", "category_id": 7, "category": {"id": 7, "name": "topology", "display_name": "Topology", "description": "Properties preserved under continuous deformations.", "slug": "topology", "order_index": 7, "created_at": "2026-07-31T15:26:25.671Z"}, "status": "open"}
% FIRST_POSTED: null
% ENTRY_KIND: statement_only
% Imported from: batch13.tex; UnsolvedMath ID 2884
\documentclass[11pt]{article}
\usepackage[margin=25mm]{geometry}
\usepackage{fontspec}
\setmainfont{Latin Modern Roman}
\usepackage{amsmath,amssymb,mathtools}
\usepackage{unicode-math}
\setmathfont{Latin Modern Math}
\usepackage{xurl,hyperref}
\hypersetup{colorlinks=true,urlcolor=blue}
\setcounter{secnumdepth}{0}
\setlength{\parindent}{0pt}
\setlength{\parskip}{5pt}
\setlength{\emergencystretch}{3em}
\newcommand{\sref}[2]{\hyperlink{source:#1:#2}{[S#2]}}
\newcommand{\cataloguescope}[1]{\paragraph{Recorded target.}#1}
\begin{document}
% UnsolvedMath ID 2884; source catalogue code KP-4.8
\setcounter{section}{2883}
\section{Kirby Problem 4.8}\label{2884}
{\small\sffamily UnsolvedMath ID 2884 · Topology}
\subsection{Definitions and mathematical statement}

\paragraph{Shared notation.}$\mathbb N=\{1,2,\ldots\}$, and $\mathbb Z,\mathbb Q,\mathbb R,\mathbb C$ denote the integers, rationals, reals and complex numbers. Any other conventions are specified locally.


Let $X$ be a closed simply connected smooth four-manifold, and let $T\subset X$ be a smoothly embedded torus with self-intersection $T\cdot T=0$ and simply connected complement $\pi_1(X\setminus T)=1$.

For a smooth knot $K\subset S^3$, let $M_K$ be zero-framed Dehn surgery on $K$, using the longitude bounding a Seifert surface in its exterior. Let $m$ be a meridian of $K$, regarded as a framed circle in $M_K$, and set $T_m=m\times S^1\subset M_K\times S^1$. If $T\subset X$ is a torus with trivial normal bundle, Fintushel--Stern knot surgery is
\[
X_K=\bigl(X\setminus\operatorname{int}(T\times D^2)\bigr)
\cup_\varphi\bigl((M_K\times S^1)\setminus\operatorname{int}(T_m\times D^2)\bigr).
\]
Here $\varphi$ is the usual orientation-reversing identification of the boundary three-tori, preserving the normal-circle class $[\{\mathrm{pt}\}\times\partial D^2]$ in the fiber-sum convention. Choose the framing and boundary identification as part of the knot-surgery data. Equivalently, the replacement piece is $S^1$ times the knot exterior, with its Seifert-longitude class identified with the normal circle of $T$. \sref{2884}{3}

A knot is prime if it is nontrivial and cannot be expressed as a connected sum of two nontrivial knots. A mirror is the image under an orientation-reversing diffeomorphism of $S^3$; isotopy of knots means ambient isotopy in $S^3$.
\paragraph{Statement recorded in the supplied source.}
For prime knots $K_1,K_2$, does
\[
 X_{K_1}\cong_{\mathrm{diff}}X_{K_2}
 \quad\Longrightarrow\quad
 K_1\text{ is isotopic to }K_2\text{ or to its mirror}
\]
hold for the above knot-surgery construction? The target is recognition of the knot up to mirroring. The hypotheses are not restricted to cusp-embedded tori, nonzero gauge invariants, or a particular starting manifold. \sref{2884}{2}
\subsection{Short English statement}
If knot surgery with two prime knots produces diffeomorphic four-manifolds, must the knots agree up to isotopy and mirror image?
\subsection{Sources}
\begin{description}
\item[\hypertarget{source:2884:1}{S1}] ULAM.AI and the UnsolvedMath Contributors, UnsolvedMath: A Curated Collection of Open Mathematics Problems, record 2884 (KP-4.8). CC BY 4.0. \url{https://huggingface.co/datasets/ulamai/UnsolvedMath}.
\item[\hypertarget{source:2884:2}{S2}] R. İnanç Baykur, Robion C. Kirby and Daniel Ruberman, K3—A New Problem List in Low-Dimensional Topology, Mathematical Surveys and Monographs 295 (2026), author version, Problem 4.8 and Remarks, PDF page 197. \url{https://math.berkeley.edu/sites/default/files/surv-295-ruberman-watermarked-author-pdf.pdf}.
\item[\hypertarget{source:2884:3}{S3}] Ronald Fintushel and Ronald J. Stern, Knots, Links, and 4-Manifolds, Inventiones Mathematicae 134 (1998), 363--400; arXiv:dg-ga/9612014v2, Introduction, pages 3--4. \url{https://arxiv.org/pdf/dg-ga/9612014v2}.
\end{description}
\label{end:2884}
\end{document}
